% ============================================================================
% 00_preliminares.tex -- Resumen / Abstract, palabras clave / keywords y
% agradecimientos. La portada (\begin{titlepage}) vive en TFT.tex porque debe
% preceder al \tableofcontents y esta acoplada al \newgeometry de nivel
% documento; este fichero solo contiene lo que va inmediatamente despues.
% ============================================================================

\begin{abstract}

  El transporte público metropolitano de Madrid, coordinado por el Consorcio
  Regional de Transportes de Madrid (CRTM), integra cuatro operadores de
  naturaleza dispar ---Metro, EMT, autobús interurbano y Cercanías---, cuya
  demanda diaria agregada condiciona decisiones operativas de flota, personal y
  frecuencias. Este trabajo evalúa si una arquitectura híbrida residual de dos
  etapas mejora la predicción de esa demanda frente a modelos consolidados de una
  sola etapa. Sobre una serie diaria continua de 1.310 observaciones (2023-01-01 a
  2026-08-02) se unifican la demanda del CRTM, la meteorología de Open-Meteo y el
  calendario laboral municipal, y se construye una matriz de 161 variables
  predictoras ---retardos de demanda total y por operador, medias y desviaciones
  móviles con desplazamiento previo estricto, armónicos de Fourier deterministas,
  codificación de calendario y meteorología retardada--- sujeta a guardias
  anti-fuga verificadas por una suite de 472 pruebas automatizadas. La primera
  etapa es una red LSTM univariante cuyas predicciones de entrenamiento se generan
  fuera del pliegue de entrenamiento (\emph{out-of-fold}) mediante una validación
  cruzada expansiva de cinco bloques; la segunda es un modelo XGBoost que corrige
  el residuo de la primera a partir de la matriz completa. El resultado central es
  negativo para la arquitectura: el híbrido no supera a un SARIMAX seleccionado por
  AIC ni a un XGBoost entrenado directamente sobre las mismas variables (MAE de
  test 234.223 frente a 163.627 y 156.121, respectivamente). Este hallazgo se
  interpreta a la luz de la literatura clásica de combinación de pronósticos
  (Granger, 1989): los conjuntos de información de ambas etapas están anidados
  ---la matriz de la etapa 2 ya representa el pasado reciente que observa la etapa
  1---, lo que deja poco margen a la complementariedad de errores que favorece a
  una combinación. Los modelos más precisos del estudio
  son dos ensambles aritméticos de SARIMAX y XGBoost ---sin ninguna red neuronal ni
  reentrenamiento---: 139.725 de MAE de test en la variante equiponderada 50/50 y
  139.681 en la ponderada por el inverso del MAE, ambas con un $R^2$ de 0,9755.
  Ambas variantes del ensamble ofrecen resultados prácticamente equivalentes. Se
  utiliza la variante equiponderada como referencia de interpretación por su
  simplicidad y por no requerir estimación de pesos, dejando constancia de que la
  preferencia entre ambas variantes se estableció a posteriori y no constituye un
  resultado confirmatorio del estudio. Su ganancia se interpreta, sin demostración
  causal, por la correlación imperfecta entre los errores de sus componentes
  (Bates y Granger, 1969). Este
  Trabajo de Fin de Máster no continúa ningún proyecto anterior propio ni ajeno, y
  toma de un trabajo
  previo de quien ejerce su dirección (Surribas-Sayago et al., 2026) únicamente inspiración de
  ingeniería de características ---la construcción de retardos, los armónicos
  deterministas y la idea de una rama exógena---, sin replicar su arquitectura: ninguna
  cifra aquí presentada procede de aquel trabajo.

  \vspace{0.2cm}

  \keywords{transporte público, predicción de demanda, series temporales
    diarias, híbrido residual LSTM-XGBoost, combinación de pronósticos,
    validación out-of-fold}
\end{abstract}

\newpage
\section*{Abstract}\label{sec:abstract}

The metropolitan public transport system of Madrid, coordinated by the Regional
Transport Consortium (CRTM), integrates four structurally different operators
---the underground network, the municipal bus company (EMT), interurban coaches,
and suburban rail (Cercan\'ias)--- whose aggregate daily demand shapes
operational decisions on fleet, staffing, and service frequency. This Master's
Thesis evaluates whether a two-stage residual hybrid architecture improves the
forecasting of that demand over consolidated single-stage models. On a continuous
daily series of 1,310 observations (1 January 2023 to 2 August 2026), the CRTM
demand data, the Open-Meteo weather data, and the municipal working calendar are
unified, and a matrix of 161 predictor variables is engineered ---lags of total
and per-operator demand, shift-first rolling means and standard deviations,
deterministic Fourier harmonics, a calendar encoding, and lagged weather---
subject to anti-leakage guards verified by a suite of 472 automated tests. The first stage is a univariate LSTM whose training-time
predictions are generated strictly out-of-fold through an expanding five-block
walk-forward cross-validation; the second is an XGBoost model that corrects the
first stage's residual using the full predictor matrix. The central result is
negative for the proposed architecture: the hybrid beats neither an
AIC-selected SARIMAX nor an XGBoost trained directly on the same variables
(test MAE of 234,223 versus 163,627 and 156,121, respectively). This finding is
interpreted in the light of the classical forecast-combination literature (Granger,
1989): the information sets of the two stages are nested ---the second stage's
feature matrix already represents the recent history the first stage observes---,
which leaves little room for the error complementarity that favours a
combination. The most accurate models in the
study are two arithmetic ensembles of SARIMAX and XGBoost ---with no neural network
and no retraining---: a test MAE of 139,725 for the 50/50 equal-weight variant and
139,681 for the inverse-MAE-weighted one, both with an $R^2$ of 0.9755. Both
variants yield practically equivalent results; the equal-weight one is used as the
interpretive reference for its simplicity and because it requires no weight
estimation, noting that the preference between the two was settled a posteriori and
is not a confirmatory result of the study. Its gain is interpreted, with no causal
demonstration, through the imperfect correlation between its components' errors
(Bates \& Granger, 1969). The contribution is therefore methodological rather than
a gain in accuracy: a negative result established under a pre-registered protocol,
with every verdict fixed in writing before the metric was observed, and interpreted
instead of being reported as an unexplained failure. This Master's Thesis continues
no earlier project
of its own or of others, and draws from a previous
work by its supervisor (Surribas-Sayago et al., 2026) only feature-engineering inspiration
---the construction of lags, the deterministic harmonics, and the idea of an exogenous
branch---, without replicating its architecture: none of the figures reported here
derives from that work.

\vspace{0.2cm}

{\small\textbf{\textsl{Keywords:\hspace{0.3cm}}} public transport, demand
forecasting, daily time series, residual LSTM-XGBoost hybrid, forecast
combination, out-of-fold validation}

\newpage
\section*{Agradecimientos}\label{sec:agradecimientos}

A la Universidad Internacional de Valencia y al profesorado del Máster
Universitario en Big Data y Ciencia de Datos, por el marco formativo en el que
se ha desarrollado este trabajo, y a la persona que ha ejercido su dirección,
por la orientación metodológica y la revisión crítica a lo largo de todo el
proceso.

Al Consorcio Regional de Transportes de Madrid y al proyecto Open-Meteo, cuyos
datos abiertos han hecho posible el estudio empírico que sustenta esta memoria.

A mi familia y a las personas cercanas, por el apoyo sostenido durante el tiempo
dedicado a este Trabajo de Fin de Máster.
